\BeleskeID{05-s055}
% element: 05-s055:d01
Znak količnika određuje se poređenjem znakova operanada, pre uzimanja apsolutnih vrednosti. Zatim se izvršava neoznačeno deljenje veličina i količnik negira ako su se znaci razlikovali. Ovo određuje zaokruživanje celobrojnog količnika prema nuli, ako neoznačeno deljenje odbacuje razlomljeni deo.

Delilac mora biti različit od nule. Najmanji negativni operand zahteva pažnju pri računanju apsolutne vrednosti, kao i pri množenju. Poseban slučaj $-2^{n-1}/(-1)=2^{n-1}$ ne može da stane u isti označeni $n$-bitni registar. Količnik i ostatak treba povezati identitetom $a=bq+r$, uz jasno odabran dogovor o znaku ostatka.
